% 2011 年北京卷文科第 16 题：甲、乙两组植树棵数的茎叶图。
% 来源 PDF：普通高考/2011/2011北京文.pdf，SHA-256 1ba0def5e75e0a4d7cc4e4e52a11a65ad598b98bc2a8e6c98df8f29fb6b9f4a6。
% 原页：第 1 页（A3 三栏）第 3 栏上部，书页码 649。
% 重绘依据：照原页——表头「甲组」「乙组」，表头下一条横线，茎列两侧各一条竖线；
%   茎 0：甲 9、9，乙 X、8、9；茎 1：甲 1、1，乙 0。乙组模糊数据以 X 表示。
\begin{tikzpicture}[semithick, scale=1.0]
  % 表头横线与茎列两条竖线
  \draw (-2.55,0.95) -- (3.05,0.95);
  \draw (-0.55,0.95) -- (-0.55,-0.55);
  \draw (0.55,0.95) -- (0.55,-0.55);

  % 表头
  \node at (-1.55,1.32) {甲组};
  \node at (1.55,1.32) {乙组};

  % 茎 0：甲 9 9，乙 X 8 9
  \node at (-1.55,0.42) {$9$};
  \node at (-1.00,0.42) {$9$};
  \node at (0,0.42) {$0$};
  \node at (1.00,0.42) {$X$};
  \node at (1.55,0.42) {$8$};
  \node at (2.10,0.42) {$9$};

  % 茎 1：甲 1 1，乙 0
  \node at (-1.55,-0.20) {$1$};
  \node at (-1.00,-0.20) {$1$};
  \node at (0,-0.20) {$1$};
  \node at (1.00,-0.20) {$0$};
\end{tikzpicture}
